% !TeX program = pdflatex
% One-file IEEE conference manuscript: references are embedded below.
% Figures are regenerated from saved results by the scripts in src/plots/.
% Compile with pdfLaTeX. Run twice when compiling locally.
% Float boundaries keep methodology floats within Methodology.
% The final float barrier keeps results before the compact bibliography.

\documentclass[conference]{IEEEtran}

\usepackage[T1]{fontenc}
\usepackage{graphicx}
\usepackage{booktabs}
\usepackage{array}
\usepackage{tabularx}
\usepackage{amsmath}
\usepackage{cite}
\usepackage{xurl}
\usepackage[section]{placeins}
\usepackage{stfloats}

% Wide figures may appear at either page margin alongside ordinary text.
% Compact, flexible spacing for ordinary single-column score tables.
\setlength{\textfloatsep}{8pt plus 2pt minus 2pt}
\setlength{\floatsep}{6pt plus 1pt minus 1pt}
\setlength{\intextsep}{6pt plus 1pt minus 1pt}
\renewcommand{\topfraction}{0.9}
\renewcommand{\bottomfraction}{0.85}
\renewcommand{\textfraction}{0.08}
\setcounter{topnumber}{3}
\setcounter{bottomnumber}{3}
\setcounter{totalnumber}{6}
\renewcommand{\dbltopfraction}{0.9}
\renewcommand{\dblfloatpagefraction}{0.8}
\setlength{\dbltextfloatsep}{8pt plus 2pt minus 2pt}
\setlength{\dblfloatsep}{6pt plus 1pt minus 1pt}
\AtBeginDocument{\raggedbottom}

\urlstyle{same}

% Ragged-right paragraph columns allow long table descriptions to wrap.
\newcolumntype{L}[1]{>{\raggedright\arraybackslash}p{#1}}
\newcolumntype{Y}{>{\raggedright\arraybackslash}X}
\newcolumntype{C}{>{\centering\arraybackslash}X}

% Shrink a natural-width table only when it exceeds the available width.
% Use inside a table/table* float around its tabular environment.
\newsavebox{\PaperTableBox}
\newenvironment{FitTable}{%
  \begin{lrbox}{\PaperTableBox}%
}{%
  \end{lrbox}%
  \ifdim\wd\PaperTableBox>\linewidth
    \resizebox{\linewidth}{!}{\usebox{\PaperTableBox}}%
  \else
    \usebox{\PaperTableBox}%
  \fi
}

% The scripts in src/plots/ regenerate these files from saved results.
% Standard graphicx search paths support both the project root and paper/.
\graphicspath{{./}%
  {figures/cross_model/}{../figures/cross_model/}%
  {figures/recalibration/}{../figures/recalibration/}%
  {figures/cross_dataset/}{../figures/cross_dataset/}%
  {figures/aggregated_pool/}{../figures/aggregated_pool/}%
  {figures/shap/catboost/}{../figures/shap/catboost/}}
% Require the actual images: missing files must not silently become blanks.
% The second argument is retained to preserve the original figure calls.
\newcommand{\PaperFigure}[2]{%
  \includegraphics[width=\linewidth,height=0.72\textheight,keepaspectratio]{#1}%
}

% Add the approved author and keywords if required.

\title{Cross-Generator Evaluation of Stylometric Classifiers for Phishing Email Authorship}

\begin{document}

\maketitle

\begin{abstract}
Large language models (LLMs) can produce fluent phishing emails,
but detectors trained on one generator may lose effectiveness
when the generator changes. Building on the benchmark of
Gutierrez et al., this study compares k-nearest neighbors (KNN),
LightGBM, CatBoost, XGBoost, and Random Forest for distinguishing
human-written from LLM-generated phishing. We reuse 9,986 emails
from five human-email sources and three generators: GPT-4.1,
DeepSeek 3.2, and Llama 3.3 70B. All 17 released stylometric
features were reproduced before reuse. Python pipelines applied
shared similarity-grouped development, calibration, and test
partitions, training-only KNN standardization, and group-aware
nested validation. Five evaluations examined intra-model,
cross-model, threshold-recalibrated, cross-dataset, and
aggregated-pool performance. LightGBM achieved the highest
mean intra-model F1 of 0.9907, whereas KNN achieved the highest
default different-generator mean F1 of 0.9120. Default transfer
gaps ranged from 5.31 to 31.90 percentage points. Recalibration
improved mean transfer F1 for every classifier, with CatBoost
reaching 0.9244; recovered misses involved false-positive
trade-offs. Pooled XGBoost achieved overall F1 of 0.9975,
with five human false positives and no generated-email misses.
Human-source holdout revealed source-dependent errors.
CatBoost SHAP reproduced the published XGBoost top-five feature
membership and identified negative mean time-pressure
contributions among missed Llama emails, providing associative
rather than causal evidence. An accountability ledger links
predictions, configurations, metrics, and hashed artifacts.
The findings demonstrate the need to assess generator transfer,
human-source variation, and threshold-dependent errors alongside
familiar-generator performance.
\end{abstract}

\section{Introduction}
\label{sec:introduction}

Phishing emails use deceptive messages to persuade recipients to
disclose information or perform actions that benefit an attacker.
Detection research has examined message content, linguistic patterns,
and machine learning as ways to identify these threats
\cite{salloum2022}. Large language models (LLMs) add to this challenge
by enabling the production of fluent, context-specific messages with
limited manual effort. Studies of generated phishing have demonstrated
both its ability to elicit recipient interaction and its capacity to
bypass some commercial email filters
\cite{heiding2024,opara2025}. These developments motivate examining
how detection methods behave when the characteristics of the
evaluated emails change.

Current research approaches this problem from several directions.
Some studies use language models to classify emails and explain
their decisions~\cite{koide2026}. Others extract stylometric features,
such as vocabulary variety, sentence structure, and persuasion
markers, and supply these measurements to conventional
classifiers~\cite{opara2025}. Research on general generated-text
detection also examines robustness across domains and generators.
MAGE demonstrated the difficulty of evaluating detectors on
unfamiliar sources~\cite{li2024mage}, while Borile and
Abrate~\cite{borile2025} investigated how dataset-specific patterns
can limit generalization. Together, these studies show why performance
on familiar data provides only part of the evidence needed to assess
a detector.

The present study examines authorship within phishing emails.
Its classifiers distinguish human-written phishing from
LLM-generated phishing. This authorship signal can support the
analysis of generated threats within a broader email-security
workflow. Evaluating it requires attention to both sides of the
classification problem. A model can miss generated emails, and it
can also incorrectly classify human-written emails as generated.

Gutierrez et al.~\cite{gutierrez2026} investigated this task using
17 stylometric features and phishing emails from three generators.
Their evaluation of Logistic Regression and XGBoost showed that
strong performance within a generator can coexist with a substantial
performance drop when the generator changes. They further examined
threshold recalibration and training on a pooled generator dataset.
Their results motivate a broader question about classifier choice.
When the same features and evaluation records are used, do different
classifiers exhibit the same transfer weaknesses, and does threshold
adjustment produce similar benefits and error trade-offs?

This question is important because a higher aggregate score can
hide changes in the types of mistakes a model makes. A lower
threshold may recover generated emails while increasing false
positives on human-written emails. Changes in the human-email
source may also affect performance, even when the generated-email
pool is held fixed. In addition, a feature that ranks highly overall
may contribute differently to missed and correctly detected emails.
Examining these effects requires connecting aggregate metrics with
individual predictions and feature contributions.

We address these questions through an independent extension of the
released benchmark. The study reuses 9,986 phishing emails and
verifies the 17 released stylometric features. It compares
k-nearest neighbors, LightGBM, CatBoost, XGBoost, and Random
Forest using shared evaluation partitions that keep identified
similar emails together. The five evaluation conditions are
intra-model testing, cross-model testing, threshold recalibration,
cross-dataset human-source holdout, and aggregated-pool training.

The study addresses four research questions:
\begin{enumerate}
    \item How does classifier choice affect performance within a
    generator and transfer to a different generator?

    \item How much does threshold recalibration improve transfer
    performance, and how does it change missed detections and
    human-email false positives?

    \item How do changes in the human-email source and pooled
    generator training affect classifier performance?

    \item Which stylometric features contribute to the examined
    predictions, and how do CatBoost's feature rankings compare
    with the published XGBoost rankings?
\end{enumerate}

The contributions are a consistent five-classifier benchmark,
an analysis of performance gaps and their associated errors,
and a focused SHAP investigation of feature contributions.
An accountability ledger records experimental settings and links
results to individual predictions and saved artifacts. This
supports verification of the reported findings and provides a
traceable foundation for future experiments.

\section{Related Work}
\label{sec:related-work}

\subsection{Phishing Email Detection}

Salloum et al.~\cite{salloum2022} examined the difficulty of detecting
phishing emails as attackers change their language and delivery
strategies. Their systematic review covered natural language processing
and machine learning approaches, including datasets, feature
representations, classification methods, and evaluation measures. The
review organized the available techniques and resources and discussed
directions for improving content analysis. It provides a foundation
for assessing how dataset selection and feature representation influence
the evidence produced by a phishing detection study.

Chataut et al.~\cite{chataut2024} investigated whether large language
models could identify phishing emails despite increasingly sophisticated
attacks. They compared GPT-3.5, GPT-4, and a customized ChatGPT
configuration using phishing and legitimate emails. Their findings
showed that the models could identify phishing attempts, with performance
varying across configurations. The study illustrates the potential of
language models as detectors and the importance of model selection.
Its classification task concerns malicious intent, while the present
study concerns human or machine authorship within phishing emails.

Koide et al.~\cite{koide2026} addressed both detection errors and the
difficulty users face in understanding why a message has been flagged.
They introduced ChatSpamDetector, which converts email content into
prompts for language-model analysis and returns a phishing decision
with supporting explanations. Their GPT-4 configuration achieved a
reported accuracy of 99.70\%. This work combines classification with
reasoning that users can inspect. Its results concern phishing detection
using a language model as the detector, providing a different setting
from feature-based authorship classification.

\subsection{Generated Phishing and Its Effectiveness}

Heiding et al.~\cite{heiding2024} investigated whether language models
could reduce the effort required to produce persuasive phishing
messages. They compared GPT-4-generated emails, manually designed
emails based on the V-Triad persuasion framework, a combination of
GPT-4 and V-Triad, and generic control emails. Their study involved
112 participants. The reported click-through range was 30--44\% for
GPT-generated emails and 19--28\% for the control group. The authors
also tested language models on identifying malicious email intent.
The findings show that generated messages can elicit user interaction,
motivating defensive evaluation. Click-through rates measure recipient
behavior and therefore require a different interpretation from the
classification metrics used in the present study.

\subsection{Stylometric Classification of Generated Text and Phishing}

Opara~\cite{opara2024styloai} examined which measurable writing
characteristics could distinguish generated text from human-authored
text. The proposed StyloAI system combined 31 stylometric features
with a Random Forest classifier and was evaluated on the
AuTextification and Education datasets. It achieved reported test
accuracies of 81\% and 98\%, respectively. These results demonstrate
that writing-style features can support authorship classification,
while the different scores illustrate the importance of evaluating
performance on each dataset. StyloAI provides a methodological
connection to the Random Forest component of the present study,
although its evaluations concern general text rather than phishing
authorship.

Opara et al.~\cite{opara2025} subsequently investigated whether
commercial filters and feature-based classifiers could identify
generated phishing. They constructed 63 phishing emails and 63
legitimate emails using GPT-4o and evaluated Gmail, Outlook, and
Yahoo filtering behavior. They also trained Logistic Regression,
a support vector machine, Random Forest, and XGBoost using 60
stylometric and linguistic features. Reported accuracies were 96\%
for XGBoost, 92\% for Random Forest and Logistic Regression, and
85\% for the support vector machine. XGBoost achieved a receiver
operating characteristic area of 0.99. The study supports examining
linguistic features with conventional classifiers, but its single-generator
phishing-versus-legitimate task differs from the authorship task
investigated here.

\subsection{Generalization Across Generators and Domains}

Li et al.~\cite{li2024mage} addressed a limitation of generated-text
detection research in which evaluation was restricted to particular
domains or language models. They introduced MAGE, a benchmark
covering ten domains and text produced by 27 language models.
The evaluation included increasingly difficult conditions involving
unseen domains and generators. Performance declined as the test
distribution differed from training. In one configuration combining
unseen domains and an unseen generator, a Longformer-based detector
achieved 84.12\% recall on generated text. These findings motivate
explicit transfer evaluations. MAGE concerns general text detection,
so its scores provide context for robustness rather than a directly
comparable phishing benchmark.

Borile and Abrate~\cite{borile2025} investigated why transformer-based
generated-text detectors generalize poorly across domains. They
examined confounding neurons that encode dataset-specific patterns
and proposed interventions that suppress their influence. Their
experiments showed reduced topic-specific bias and improved
generalization across domains. This work highlights how a detector
can depend on patterns associated with its training data. It motivates
examining source changes and model behavior in the present study,
while explanations of individual predictions remain evidence about
the fitted classifier's behavior.

\subsection{Cross-Generator Phishing Authorship Detection}

Gutierrez et al.~\cite{gutierrez2026} directly investigated whether
stylometric phishing authorship detectors transfer between generators.
They assembled 9,986 emails comprising human-written phishing
and phishing generated by GPT-4.1, DeepSeek 3.2, and Llama 3.3 70B.
Using 17 stylometric features, they evaluated Logistic Regression
and XGBoost under intra-model, cross-model, threshold-recalibrated,
cross-dataset, and aggregated-pool conditions.

XGBoost achieved intra-generator F1 scores from 0.9516 to 0.9650.
The reported mean same-generator versus different-generator F1 gap
was 28.0 percentage points at the default threshold. Although F1
declined under generator changes, receiver operating characteristic
area remained above 0.96 in every different-generator cell. Threshold
recalibration reduced the gap to 4.0 percentage points, while pooled
training achieved a reported F1 of 0.997 for each generator. Their
SHAP analysis identified politeness density and type-token ratio
among the five most important features for all three generator-specific
XGBoost models. These findings established the benchmark and
evaluation questions extended by the present study.

The present research compares k-nearest neighbors, LightGBM,
CatBoost, XGBoost, and Random Forest using the released features
and consistent evaluation partitions. It extends the analysis through
group-aware partitioning, prediction-level accountability records,
and examination of missed detections and human-email false
positives. A focused SHAP comparison investigates feature
contributions on the same Llama emails. Comparisons with published
scores account for differences in class definitions, data partitions,
model selection, and metric aggregation.

\section{Methodology}
\label{sec:methodology}

\subsection{Research Design and Dataset}

This study extended the phishing authorship benchmark released by
Gutierrez et al.~\cite{gutierrez2026}. We reused the released dataset
and its 17 stylometric features to compare five classifiers under
consistent experimental conditions. The classification target was
authorship within phishing emails. Human-written phishing was
assigned label 0, and language-model-generated phishing was assigned
label 1.

The corpus contained 9,986 emails, comprising 5,000 human-written
emails and 4,986 generated emails. The generated subset contained
1,665 GPT-4.1 emails, 1,665 DeepSeek 3.2 emails, and 1,656
Llama 3.3 70B emails. The human subset combined CEAS-08,
TREC-07, Nazario, Nigerian Fraud, and a fraud-labeled Enron subset.
The released source and authorship labels were retained. The
generated emails covered banking, parcel delivery, IT support,
tax, and human resources. The dataset and reference code are
available in the authors' version-1 Zenodo
archive~\cite{gutierrezdataset2026}; this study reused the emails
and did not generate a new corpus.

Original dataset files were preserved. File hashes, record
identifiers, source labels, and feature order were checked to
maintain traceability between the released data and the experiments.

\subsection{Stylometric Features and Validation}

All five classifiers received the same 17 numerical features,
organized into lexical, syntactic, stylistic, and phishing-specific
categories. Table~\ref{tab:stylometric-features} describes their
calculations. Features were reproduced from the stored email text
using the released extraction conventions, including the original
lexicons and rounding precision. The stored text combined subject
and body content and retained its original formatting.

Reproduction used spaCy 3.8.16 with the English
\texttt{en\_core\_web\_sm} model, version 3.8.0. The full-corpus
validation reproduced all 17 stored feature values for all 9,986
records after rounding. This established agreement with the released
implementation. The subsequent analyses examined the behavior
of those features within the fitted classifiers.

\begin{table*}[!tp]
\centering
\small
\caption{The 17 reused stylometric features and their extraction
conventions. Here, $N$ is the number of alphabetic word tokens,
$T$ is the number of non-whitespace tokens, and $S$ is the number
of sentences.}
\label{tab:stylometric-features}
\renewcommand{\arraystretch}{1.12}
\begin{tabularx}{\textwidth}{
    @{}L{0.14\textwidth}L{0.25\textwidth}Y@{}
}
\toprule
Category & Feature & Measurement \\
\midrule

Lexical
& Type-token ratio
& Number of distinct lowercased words divided by $N$. \\

& Mean word length
& Mean number of characters per alphabetic word token. \\

& Mean sentence length
& $T/S$, including punctuation tokens in $T$. \\

& Yule's $K$
& Word-frequency concentration, calculated as
$10^{4}(\sum_w f_w^2-N)/N^2$, where $f_w$ is the frequency
of word $w$. \\

\midrule

Syntactic
& Clause density
& Number of tokens with dependency labels
\texttt{ROOT}, \texttt{ccomp}, \texttt{advcl},
\texttt{relcl}, or \texttt{xcomp}, divided by $S$. \\

& Noun ratio
& Number of tokens tagged \texttt{NOUN}, divided by $T$. \\

& Verb ratio
& Number of tokens tagged \texttt{VERB}, divided by $T$. \\

& Mean dependency-parse depth
& Mean token distance from the dependency root. The released
convention includes whitespace tokens and caps depth at 200. \\

\midrule

Stylistic
& Imperative count
& Heuristic count of sentences whose first non-whitespace token
has tag \texttt{VB} and dependency label \texttt{ROOT}
or \texttt{ccomp}. \\

& First-person pronoun ratio
& Matches to the released first-person pronoun list, divided by $N$. \\

& Second-person pronoun ratio
& Matches to the released second-person pronoun list, divided by $N$. \\

& Politeness density
& Matches to the released politeness lexicon per 100 word tokens. \\

& Urgency density
& Matches to the released urgency lexicon per 100 word tokens. \\

\midrule

Phishing-specific
& URL density
& URL pattern matches per 100 word tokens. \\

& Call-to-action density
& Matches to the released action-request lexicon per 100 word tokens. \\

& Authority-appeal density
& Matches to the released authority lexicon per 100 word tokens. \\

& Time-pressure density
& Matches to the released deadline and time-pressure lexicon
per 100 word tokens. \\

\bottomrule
\end{tabularx}
\end{table*}

\subsection{Partitioning and Similarity Groups}

To reduce overlap between partitions, emails were grouped using
text similarity. For grouping, whitespace was collapsed and text
was case-folded. Each email was represented by binary word
trigrams, and pairs with cosine similarity of at least 0.88 were
linked. Connected components formed the final similarity groups.
This representation was used exclusively for grouping.

Group-aware stratification approximately balanced authorship,
source, and generated-email category. The resulting allocation
contained 5,991 development records, 1,998 calibration records,
and 1,997 final-test records. Development contained 3,000 human
emails and 2,991 generated emails. Calibration contained 1,000
human emails and 998 generated emails, while testing contained
1,000 human emails and 997 generated emails.

Record and group identifiers were saved in split manifests.
Checks confirmed that no identified similarity group crossed
these partitions. Within development, validation folds also kept
each group together. Whitespace normalization used for grouping
did not replace the stored text or feature values.

\subsection{Classifiers and Model Selection}

We evaluated k-nearest neighbors (KNN), LightGBM, CatBoost,
XGBoost, and Random Forest. Each classifier used all 17 features.
KNN included a standardization step fitted within each training
fold. The tree classifiers used the original numerical feature
scales. Table~\ref{tab:classifier-settings} explains how the
classifiers used the features and lists the parameter searches.

Each classifier received a grid of 20 parameter combinations.
Selection maximized development F1 at a decision threshold of
0.5 using shared group-aware folds. Randomized operations used
seed 42 unless a fold-specific seed was required. Computation
used CPU execution and one worker per fitting operation.
Experiments were implemented in Python 3.13.5 using
scikit-learn 1.9.1, LightGBM 4.6.0, CatBoost 1.2.10,
and XGBoost 3.4.1. Configurations, package versions, and
selected parameters were recorded for each run.

\begin{table*}[!tp]
\centering
\small
\caption{Classifier use of the shared features and parameter-search
spaces. Every classifier evaluated 20 parameter combinations.}
\label{tab:classifier-settings}
\renewcommand{\arraystretch}{1.15}
\begin{tabularx}{\textwidth}{
    @{}L{0.13\textwidth}L{0.33\textwidth}Y@{}
}
\toprule
Classifier & How the features influence predictions & Search space \\
\midrule

KNN
&
Distances between standardized feature vectors determine neighboring
training emails and their votes.
&
Neighbors $\{3,5,11,21,31\}$; weights
$\{\text{uniform},\text{distance}\}$; Minkowski parameter
$\{1,2\}$.
\\

LightGBM
&
Boosted trees learn feature thresholds and combinations that
change the generated-email score.
&
Trees $\{100,200\}$; leaves $\{7,15,31,63,127\}$;
learning rate $\{0.05,0.1\}$.
\\

CatBoost
&
Symmetric boosted trees combine feature splits to produce
the classification score.
&
Iterations $\{100,200\}$; depth $\{3,4,5,6,7\}$;
learning rate $\{0.05,0.1\}$.
\\

XGBoost
&
Boosted trees combine feature thresholds while fitting the
binary classification objective.
&
Trees $\{100,200\}$; depth $\{2,3,4,6,8\}$;
learning rate $\{0.05,0.1\}$.
\\

Random Forest
&
Trees trained on bootstrap samples use feature splits,
and their class-probability estimates are averaged.
&
Trees $\{100,200\}$; maximum depth
$\{4,8,12,16,\text{unrestricted}\}$; minimum samples per
leaf $\{1,4\}$. The feature subset size was fixed to
\texttt{sqrt}.
\\

\bottomrule
\end{tabularx}
\end{table*}

\subsection{Evaluation Conditions}

\subsubsection{Intra-Model Evaluation}

For each generator, a classifier distinguished human-written
phishing from phishing produced by that generator. Evaluation
used development records in nested, group-aware cross-validation.
Five outer folds measured performance. Within each outer
training fold, five inner folds selected among the 20 parameter
settings. The selected pipeline was then refitted on that outer
training fold and evaluated on its excluded outer fold.

This procedure required 505 fits per classifier and generator,
including the five outer-training refits. Performance was
summarized using the mean and sample standard deviation across
outer folds. The standard deviation describes fold variation.
The outer training sets overlap, so these five scores were
not treated as independent samples for significance testing.

\subsubsection{Cross-Model Evaluation}

Separate development models were selected and refitted using
human emails and one generator. Each fitted model was evaluated
on held-out human emails and emails from each of the three
generators. This produced a $3 \times 3$ matrix per classifier,
with three same-generator and six different-generator cells.

All classifiers used the same target-specific test records.
A generated-email score of at least 0.5 assigned label 1.
The same held-out human emails were reused across matrix
cells to support comparisons under changes in generator.

\subsubsection{Threshold Recalibration}

The fitted cross-model classifiers were retained while their
decision thresholds were selected using separate calibration
emails. Each training-generator and target-generator pair used
the target generator's calibration emails together with the
human calibration subset.

Candidate thresholds consisted of the distinct calibration
scores and 0.5. The threshold maximizing calibration F1 was
selected. Ties favored the threshold closest to 0.5, followed
by the larger threshold. The chosen threshold was applied to
the original test scores. Here, recalibration means adjustment
of the decision cutoff; it does not transform the scores or
establish probability calibration. We recorded F1 changes,
recovered generated-email misses, newly missed emails,
additional human false positives, and cleared human false positives.

\subsubsection{Cross-Dataset Evaluation}

To examine changes in the human-email source, we held out
each human corpus in turn. Models were trained and tuned using
development emails from the other four human sources and
the fixed development pool of all three generators. Source
exclusion also respected similarity groups.

Each fitted model was evaluated on the reserved test records.
For each human source, its false-positive rate when excluded
was compared with the mean rate from the four models that
included that source. These comparisons were descriptive,
because excluding different sources changed the size and
composition of the human training subset.

\subsubsection{Aggregated-Pool Evaluation}

One model per classifier was trained using the full development
partition, combining human emails with emails from all three
generators. Parameter selection used five group-aware development
folds, followed by refitting on the complete development set.

Each model scored the 1,997 test emails once. We reported
overall pooled metrics and separate generator-specific metrics.
Each generator-specific calculation combined that generator's
test emails with the same 1,000 human test emails. Pooled
training therefore evaluated new records from generators
represented in training.

\subsection{Performance Measures and Transfer Gap}

Evaluation included accuracy, precision, recall, F1, sensitivity,
specificity, receiver operating characteristic area, average
precision, and confusion matrices. The positive class was
generated phishing. A false positive was a human-written email
classified as generated. A false negative was a generated email
classified as human-written.

For classifier $c$, the transfer gap was calculated from its
held-out cross-model matrix:
\begin{equation}
\Delta_c =
\frac{1}{3}\sum_{g=1}^{3} F_{1,c}^{g,g}
-
\frac{1}{6}
\sum_{\substack{g,h=1\\g\ne h}}^{3} F_{1,c}^{g,h}.
\label{eq:transfer-gap}
\end{equation}
Here, $F_{1,c}^{g,h}$ is the F1 of classifier $c$ trained on
generator $g$ and tested on generator $h$. Multiplying
$\Delta_c$ by 100 expresses the gap in percentage points.
The same definition was used before and after recalibration.

Human-source false-positive rates were examined for all five
classifiers. Additional theme-level threshold-error summaries
were available for KNN, LightGBM, and CatBoost. Counts were
reported with subgroup denominators. Comparisons retained
email identifiers to account for paired predictions.

\subsection{SHAP Analysis and Comparison with Published XGBoost}

Shapley additive explanations (SHAP) assign feature contributions
to individual model outputs~\cite{lundberg2017}. We used
CatBoost's native SHAP implementation to explain its raw
classification scores. Each email in the intra-model analysis
was explained by the outer-fold model that excluded it during
training. Contributions were summarized across all outer-fold
evaluation records, including human emails and emails from
the relevant generator.

Global importance was calculated as the mean absolute SHAP
contribution for each feature. We compared the membership
and ordering of CatBoost's five highest-ranked features with
the published XGBoost rankings of Gutierrez
et al.~\cite{gutierrez2026}. Interpretation accounted for
differences in models, training procedures, and explanation
datasets.

A focused diagnostic examined DeepSeek-trained models on
Llama test emails. CatBoost explanations were compared between
missed and detected Llama emails. We additionally computed
LightGBM feature contributions on the same records and examined
the subsets both models missed, both detected, or classified
differently. For each explanation, the base score plus feature
contributions was checked against the model's raw output,
and the prediction score was checked against the saved result.

Positive contributions increased the score for generated
authorship, while negative contributions reduced it. These
analyses describe associations with model predictions and
provide evidence for investigating errors. Causal explanations
would require additional controlled experiments.

\subsection{Accountability Ledger}

An accountability ledger connected prediction records with
experiment-level reports. Prediction records stored the
experiment identifier, email identifier, similarity group,
actual label, predicted label, score, and the applicable
evaluation context. Depending on the evaluation, this context
included the training generator, test generator, fold, or
excluded human source. Explanation records linked SHAP
contributions to the corresponding email and saved model.

Experiment-level JSON reports stored dataset and split
identifiers, feature order, classifier configuration, selected
parameters, evaluation metrics, and relevant performance-gap
or recalibration results. SHA-256 hashes identified input
files and saved artifacts. Software versions and random seeds
supported reproducibility.

The ledger provides a record for later verification and further
research. Future experiments can extend these records while
retaining the identifiers needed to trace aggregate findings
back to individual predictions.


\section{Results}
\label{sec:results}

\subsection{Intra-Model Performance}

Table~\ref{tab:intra-results} reports outer-fold F1 means and
sample standard deviations. LightGBM had the highest mean
across generators (0.9907). XGBoost led on GPT-4.1, while
LightGBM led on DeepSeek and Llama. These are descriptive
fold results, not significance-tested differences.

\begin{table}[!htbp]
\centering
\setlength{\tabcolsep}{2pt}
\renewcommand{\arraystretch}{1.05}
\footnotesize
\caption{Intra-model F1 under nested group-aware cross-validation.
Each generator entry gives the mean and sample standard deviation
across five outer folds. The final column averages the three
generator means equally.}
\label{tab:intra-results}
\begin{FitTable}
\begin{tabular}{@{}lcccc@{}}
\toprule
Classifier & GPT-4.1 & \shortstack{DeepSeek\\3.2} & \shortstack{Llama\\3.3 70B}
& \shortstack{Mean\\F1} \\
\midrule
KNN
& \shortstack{$0.9503$\\$\pm 0.0172$}
& \shortstack{$0.9550$\\$\pm 0.0119$}
& \shortstack{$0.9450$\\$\pm 0.0097$}
& 0.9501 \\
LightGBM
& \shortstack{$0.9895$\\$\pm 0.0081$}
& \shortstack{$0.9910$\\$\pm 0.0057$}
& \shortstack{$0.9915$\\$\pm 0.0058$}
& 0.9907 \\
CatBoost
& \shortstack{$0.9896$\\$\pm 0.0061$}
& \shortstack{$0.9901$\\$\pm 0.0072$}
& \shortstack{$0.9890$\\$\pm 0.0042$}
& 0.9895 \\
XGBoost
& \shortstack{$0.9900$\\$\pm 0.0063$}
& \shortstack{$0.9886$\\$\pm 0.0062$}
& \shortstack{$0.9869$\\$\pm 0.0062$}
& 0.9885 \\
Random Forest
& \shortstack{$0.9850$\\$\pm 0.0078$}
& \shortstack{$0.9869$\\$\pm 0.0075$}
& \shortstack{$0.9823$\\$\pm 0.0080$}
& 0.9847 \\
\bottomrule
\end{tabular}
\end{FitTable}
\end{table}

\subsection{Cross-Model Performance and Degradation}

KNN had the highest mean different-generator F1 at the default
threshold (0.9120), despite the lowest mean intra-model F1.
Default transfer gaps ranged from 5.31 to 31.90 percentage
points (Table~\ref{tab:transfer-results}). DeepSeek-to-Llama
was the weakest pair for every classifier
(Fig.~\ref{fig:cross-model-results}). Gaps use the diagonal
and off-diagonal cells of the held-out matrix; the nested
cross-validation scores are separate estimates.

\begin{table}[!htbp]
\centering
\setlength{\tabcolsep}{1.5pt}
\renewcommand{\arraystretch}{1.05}
\footnotesize
\caption{Held-out cross-model F1 and transfer gaps.
Same-generator F1 averages the three diagonal cells.
Different-generator F1 averages the six off-diagonal cells.
Gaps are expressed in percentage points (pp).}
\label{tab:transfer-results}
\begin{tabularx}{\columnwidth}{@{}L{0.20\columnwidth}CCCCC@{}}
\toprule
Classifier & KNN & LightGBM & CatBoost & XGBoost & \shortstack{Random\\Forest} \\
\midrule
\shortstack[l]{Default\\same-\\generator\\F1} & 0.9652 & 0.9945 & 0.9945 & 0.9930 & 0.9904 \\
\shortstack[l]{Default\\different-\\generator\\F1} & 0.9120 & 0.7330 & 0.7556 & 0.7182 & 0.6714 \\
\shortstack[l]{Recalibrated\\different-\\generator\\F1} & 0.9213 & 0.9153 & 0.9244 & 0.8962 & 0.9027 \\
\shortstack[l]{Default\\gap (pp)} & 5.31 & 26.15 & 23.89 & 27.48 & 31.90 \\
\shortstack[l]{Recalibrated\\gap (pp)} & 6.24 & 7.67 & 6.96 & 9.57 & 8.92 \\
\bottomrule
\end{tabularx}
\end{table}

\begin{figure*}[!tbp]
\centering
\PaperFigure
  {cross_model_f1_before_after_recalibration.png}
  {0.65\textheight}
\caption{Cross-model F1 at the default and recalibrated thresholds.
Rows within each panel identify the training generator and columns
identify the test generator. All panels use the same color scale.
Recalibration uses separate calibration records and retains the
original held-out test scores.}
\label{fig:cross-model-results}
\end{figure*}

\subsection{Threshold Recalibration and Error Trade-Offs}

Recalibration increased mean different-generator F1 for every
classifier; CatBoost reached the highest mean (0.9244).
KNN's gap nevertheless widened because its same-generator
mean increased more. Individual cells could also decline:
KNN's Llama-to-DeepSeek F1 changed from 0.9066 to 0.9014.

For DeepSeek-to-Llama, all selected cutoffs were below 0.5.
Recovered generated-email misses were accompanied by additional
human false positives (Table~\ref{tab:threshold-errors} and
Fig.~\ref{fig:threshold-tradeoff}). Scores were unchanged,
so AUROC and AP were unchanged.

\begin{table}[!htbp]
\centering
\setlength{\tabcolsep}{1.5pt}
\renewcommand{\arraystretch}{1.05}
\footnotesize
\caption{DeepSeek-to-Llama threshold recalibration on the same
329 Llama and 1,000 human test emails. FN denotes generated
emails classified as human-written. FP denotes human emails
classified as generated. Arrows indicate default to recalibrated
results. AUROC is the area under the receiver operating
characteristic curve, and AP is average precision.}
\label{tab:threshold-errors}
\begin{tabularx}{\columnwidth}{@{}L{0.20\columnwidth}CCCCC@{}}
\toprule
Classifier & KNN & LightGBM & CatBoost & XGBoost & \shortstack{Random\\Forest} \\
\midrule
\shortstack[l]{Selected\\threshold} & 0.291586 & \shortstack{$1.305857$\\$\times 10^{-5}$} & 0.004463 & 0.001524 & 0.089987 \\
F1 & \shortstack{$0.8298$\\$\to 0.8709$} & \shortstack{$0.4977$\\$\to 0.8513$} & \shortstack{$0.5577$\\$\to 0.8464$} & \shortstack{$0.4771$\\$\to 0.8042$} & \shortstack{$0.4313$\\$\to 0.8304$} \\
FN & \shortstack{$78$\\$\to 39$} & \shortstack{$219$\\$\to 60$} & \shortstack{$201$\\$\to 48$} & \shortstack{$225$\\$\to 60$} & \shortstack{$238$\\$\to 50$} \\
FP & \shortstack{$25$\\$\to 47$} & \shortstack{$3$\\$\to 34$} & \shortstack{$2$\\$\to 54$} & \shortstack{$3$\\$\to 71$} & \shortstack{$2$\\$\to 64$} \\
AUROC / AP & \shortstack{0.9264\\ / 0.8471} & \shortstack{0.9733\\ / 0.9292} & \shortstack{0.9740\\ / 0.9317} & \shortstack{0.9566\\ / 0.8939} & \shortstack{0.9705\\ / 0.9165} \\
\bottomrule
\end{tabularx}
\end{table}

\begin{figure*}[!tbp]
\centering
\PaperFigure
  {deepseek_to_llama_threshold_tradeoff.png}
  {0.34\textheight}
\caption{Changes in missed Llama emails and human-email false
positives after recalibration for DeepSeek-trained classifiers.
Open circles indicate default-threshold counts and diamonds
indicate recalibrated counts.}
\label{fig:threshold-tradeoff}
\end{figure*}

\subsection{Cross-Dataset Human-Source Evaluation}

Nazario holdout produced the largest false-positive-rate increase
for every classifier; held-out rates ranged from 20.5\% to
50.5\% (Table~\ref{tab:nazario-results}). All five classifiers
produced zero false positives on held-out CEAS-08 emails.
Figure~\ref{fig:source-results} shows the source-specific
variation. Excluding different human corpora also changed
training-set size and composition.

\begin{table}[!htbp]
\centering
\small
\caption{False positives on the 200 Nazario test emails.
The included-source rate is averaged across the four models
whose training sets contained Nazario.}
\label{tab:nazario-results}
\setlength{\tabcolsep}{4pt}
\begin{FitTable}
\begin{tabular}{@{}lccc@{}}
\toprule
Classifier
& \shortstack{Included\\FPR (\%)}
& \shortstack{Held-out\\FPR (\%)}
& \shortstack{Held-out\\FP count} \\
\midrule
KNN           & 14.875 & 50.5 & 101 \\
LightGBM      &  2.500 & 21.5 &  43 \\
CatBoost      &  3.250 & 26.5 &  53 \\
XGBoost       &  2.625 & 22.5 &  45 \\
Random Forest &  1.625 & 20.5 &  41 \\
\bottomrule
\end{tabular}
\end{FitTable}
\end{table}

\begin{figure*}[!tbp]
\centering
\PaperFigure
  {human_source_false_positive_rates.png}
  {0.45\textheight}
\caption{Human-email false-positive rates when each source
was included in training or held out. Included-source rates
average the four applicable models. Each comparison uses
the same held-out emails from the source being examined.}
\label{fig:source-results}
\end{figure*}

\subsection{Aggregated-Pool Performance}

XGBoost had the highest overall pooled F1 (0.9975), with
five human false positives and no generated-email misses.
Its advantage over LightGBM amounted to one human-email
prediction on this test set. Table~\ref{tab:pooled-results}
reports all metrics and error counts.

\begin{table}[!htbp]
\centering
\setlength{\tabcolsep}{1.5pt}
\renewcommand{\arraystretch}{1.05}
\footnotesize
\caption{Aggregated-pool results on 1,000 human and 997 generated
test emails. Recall is also the sensitivity of the generated-email
class. FP and FN are counts. AUROC and AP use the continuous
prediction scores and are shown to six decimal places.}
\label{tab:pooled-results}
\begin{tabularx}{\columnwidth}{@{}L{0.20\columnwidth}CCCCC@{}}
\toprule
Classifier & KNN & LightGBM & CatBoost & XGBoost & \shortstack{Random\\Forest} \\
\midrule
Accuracy & 0.9805 & 0.9970 & 0.9965 & 0.9975 & 0.9950 \\
Precision & 0.9624 & 0.9940 & 0.9940 & 0.9950 & 0.9960 \\
Recall & 1.0000 & 1.0000 & 0.9990 & 1.0000 & 0.9940 \\
Specificity & 0.9610 & 0.9940 & 0.9940 & 0.9950 & 0.9960 \\
F1 & 0.9808 & 0.9970 & 0.9965 & 0.9975 & 0.9950 \\
AUROC & 0.990957 & 0.999948 & 0.999945 & 0.999932 & 0.999763 \\
AP & 0.982198 & 0.999947 & 0.999944 & 0.999931 & 0.999743 \\
FP & 39 & 6 & 6 & 5 & 4 \\
FN & 0 & 0 & 1 & 0 & 6 \\
\bottomrule
\end{tabularx}
\end{table}

\subsubsection{Pooled Versus Individual-Generator Training}

Pooling increased Llama F1 relative to the mean of the two
models trained on other generators, by approximately
8.8--49.9 percentage points
(Fig.~\ref{fig:pool-change}). Against matched-generator
training, however, pooling reduced KNN and LightGBM F1 on
all three generators (Table~\ref{tab:matched-pool}).

The highest matched-generator F1 was 0.9985 for GPT-4.1
and 0.9940 for DeepSeek, both with CatBoost, and 0.9955 for
Llama with LightGBM. Overall pooled metrics use 997 generated
and 1,000 human emails; generator-specific metrics use
approximately 330 generated and the same 1,000 human emails.

\begin{table}[!htbp]
\centering
\setlength{\tabcolsep}{2pt}
\renewcommand{\arraystretch}{1.05}
\footnotesize
\caption{Default-threshold F1 for matched individual-generator
training and pooled training. Each pair uses identical target
test records. Matched models were trained using human emails
and the indicated generator.}
\label{tab:matched-pool}
\begin{FitTable}
\begin{tabular}{@{}lcccccc@{}}
\toprule
& \multicolumn{2}{c}{GPT-4.1}
& \multicolumn{2}{c}{\shortstack{DeepSeek\\3.2}}
& \multicolumn{2}{c}{\shortstack{Llama\\3.3 70B}} \\
\cmidrule(lr){2-3}
\cmidrule(lr){4-5}
\cmidrule(l){6-7}
Classifier & Matched & Pooled & Matched & Pooled
& Matched & Pooled \\
\midrule
KNN           & 0.9653 & 0.9448 & 0.9639 & 0.9448 & 0.9662 & 0.9440 \\
LightGBM      & 0.9955 & 0.9911 & 0.9925 & 0.9911 & 0.9955 & 0.9910 \\
CatBoost      & 0.9985 & 0.9911 & 0.9940 & 0.9896 & 0.9909 & 0.9910 \\
XGBoost       & 0.9970 & 0.9926 & 0.9910 & 0.9926 & 0.9909 & 0.9925 \\
Random Forest & 0.9925 & 0.9940 & 0.9880 & 0.9895 & 0.9909 & 0.9894 \\
\bottomrule
\end{tabular}
\end{FitTable}
\end{table}

\begin{figure*}[!tbp]
\centering
\PaperFigure
  {pooled_training_f1_change.png}
  {0.33\textheight}
\caption{Change in generator-specific F1 from pooled training
relative to the mean of the two models trained on other
generators. Positive values indicate higher pooled F1.
The pooled models have access to all three training generators
and more total training records than an individual-generator
model.}
\label{fig:pool-change}
\end{figure*}

\subsection{SHAP Comparison and Error Investigation}

\subsubsection{Published XGBoost and CatBoost Rankings}

CatBoost and the published XGBoost analysis had identical
top-five feature membership for each generator
(Table~\ref{tab:shap-rank-comparison}). Ordering was unchanged
for GPT-4.1; positions 3 and 4 exchanged for DeepSeek,
and positions 1 and 2 exchanged for Llama.

\begin{table*}[!tbp]
\centering
\small
\caption{Top-five features in descending order of mean absolute
SHAP contribution. The XGBoost column reproduces the ranking
reported in the reference paper. TTR denotes type-token ratio.}
\label{tab:shap-rank-comparison}
\renewcommand{\arraystretch}{1.15}
\begin{tabularx}{\textwidth}{
    @{}L{0.12\textwidth}YYL{0.17\textwidth}@{}
}
\toprule
Generator & Published XGBoost & Our CatBoost
& Observed ordering difference \\
\midrule

GPT-4.1
&
Politeness; mean word length; call-to-action density;
time-pressure density; TTR.
&
Politeness; mean word length; call-to-action density;
time-pressure density; TTR.
&
Same order.
\\

DeepSeek 3.2
&
Mean word length; time-pressure density; politeness;
call-to-action density; TTR.
&
Mean word length; time-pressure density; call-to-action
density; politeness; TTR.
&
Positions 3 and 4 exchanged.
\\

Llama 3.3 70B
&
Politeness; TTR; verb ratio; mean sentence length;
first-person pronoun ratio.
&
TTR; politeness; verb ratio; mean sentence length;
first-person pronoun ratio.
&
Positions 1 and 2 exchanged.
\\

\bottomrule
\end{tabularx}
\end{table*}

\subsubsection{Feature Contributions on Missed Llama Emails}

DeepSeek-trained CatBoost assigned time-pressure density a
mean signed SHAP contribution of $-0.593$ on 201 missed
Llama emails and $+0.452$ on 128 detected emails
(Fig.~\ref{fig:shap-errors}). Two-thirds of the misses had
negative time-pressure contributions. Politeness still had
a positive mean contribution among misses, illustrating
the combined influence of multiple features.

LightGBM and CatBoost both missed 187 Llama emails, both
detected 96, and disagreed on 46. For both models, mean
time-pressure contributions were negative on shared misses
and positive on shared detections. These are exploratory
associations with model scores.

\begin{figure*}[!tbp]
\centering
\PaperFigure
  {catboost_shap_llama_missed_vs_detected.png}
  {0.34\textheight}
\caption{Mean signed SHAP contributions for DeepSeek-trained
CatBoost on Llama test emails. The figure shows the three
features with the most negative mean contribution among
the 201 missed emails and compares them with the 128
detected emails. Negative contributions move the score
toward human-written authorship; positive contributions
move it toward generated authorship.}
\label{fig:shap-errors}
\end{figure*}

\section{Discussion}
\label{sec:discussion}

\subsection{Interpretation of the Findings}

Classifier choice affected transfer more than the intra-model
ranking suggested. LightGBM led the nested evaluation,
whereas KNN transferred best at the default threshold and
CatBoost had the highest recalibrated transfer mean.
Selection based only on familiar-generator performance
would therefore overlook an important weakness.

Threshold adjustment recovered many misses but increased human
false positives. For DeepSeek-to-Llama, LightGBM had AUROC
0.9733 while its default F1 was 0.4977; adjusting the cutoff
raised F1 to 0.8513. This supports examining ranking quality
and decision errors separately. It does not establish that
the scores are calibrated probabilities.

Pooling performed strongly when all three generators were
represented in training, but did not consistently improve
matched-generator F1. The Nazario results also show that
generator diversity does not remove sensitivity to the human
corpus. Both authorship classes require source-aware evaluation.

\subsection{Comparison with Previous Work}

Table~\ref{tab:prior-work-comparison} separates the classification
tasks and metrics. Accuracy on phishing-versus-legitimate email
is not directly comparable with authorship F1. For the closest
numerical comparison, the table uses our XGBoost results;
the five-classifier comparisons appear in
Tables~\ref{tab:intra-results}--\ref{tab:matched-pool}.

\begin{table*}[!tp]
\centering
\small
\caption{Quantitative comparison with relevant prior work on
phishing detection and generated-phishing authorship.
Published scores retain their original metrics. All pooled
F1 values in the two authorship rows are generator-specific,
not overall pooled scores. Different tasks and protocols
prevent a common performance ranking.}
\label{tab:prior-work-comparison}
\renewcommand{\arraystretch}{1.12}
\begin{tabularx}{\textwidth}{
  @{}L{0.13\textwidth}L{0.21\textwidth}L{0.16\textwidth}
  YL{0.16\textwidth}@{}
}
\toprule
Study & Task and data & Approach
& Reported performance & Generator-transfer assessment \\
\midrule
Opara et al.~\cite{opara2025}
& Phishing versus legitimate; 63 emails per class generated
with GPT-4o.
& 60 stylometric and linguistic features; XGBoost.
& Accuracy: 96\%; AUROC: 0.99.
& Single-generator classification; no cross-generator
authorship score reported. \\
\midrule
Koide et al.~\cite{koide2026}
& Phishing versus legitimate; 1,010 phishing and 1,000
legitimate emails from public sources.
& ChatSpamDetector with GPT-4 as the detector.
& Accuracy: 99.70\%.
& GPT-4 is a detector, not the generator of an
authorship benchmark. \\
\midrule
Gutierrez et al.~\cite{gutierrez2026}
& Human versus generated phishing; 9,986 records;
GPT-4.1, DeepSeek 3.2, and Llama 3.3 70B.
& 17 features; Logistic Regression and XGBoost;
scores shown for XGBoost.
& Intra-CV F1: 0.9516--0.9650.
Pooled F1: 0.997 for each generator.
& Nine generator pairs; default to recalibrated
gap: 28.0 to 4.0 pp. \\
\midrule
Present study
& Same released corpus; shared development, calibration,
and test partitions with similarity groups kept together.
& 17 features; five classifiers;
XGBoost scores shown for the closest-family comparison.
& XGBoost intra-CV F1: 0.9869--0.9900.
Pooled F1: 0.9926 (GPT), 0.9926 (DeepSeek),
0.9925 (Llama).
& Nine pairs per classifier; XGBoost gap:
27.48 to 9.57 pp. \\
\bottomrule
\end{tabularx}
\end{table*}

Our XGBoost gap was numerically close to the published default
gap, but remained larger after recalibration.
The partitioning, group handling, parameter search,
calibration allocation, and evaluation records differed.
These results extend the benchmark under a different protocol
and do not establish superiority over the published study.

\subsection{Feature Interpretation and Accountability}

The shared top-five features support agreement with the
published importance pattern. Time pressure was already
important in that study; our additional finding concerns
its signed contribution on missed versus detected Llama
emails. This distinction adds an error-focused diagnostic,
not evidence that time pressure causes transfer failure.

The ledger connects these observations to saved scores,
thresholds, identifiers, and hashes. It supports later
verification and investigation. Practical use would also
require an evaluation containing legitimate emails and
explicit costs for each error type.

\subsection{Limitations}

The study reused one released corpus, three generators,
five generated-email themes, and inherited source labels.
It used English lexicons and an English spaCy model;
language labels were not independently audited, and no
multilingual evaluation was conducted. Conclusions therefore
concern this English-oriented benchmark.

Feature reproduction establishes agreement with the released
implementation, rather than complete linguistic validity.
Formatting was retained: the diagnostic audit found spaCy
whitespace tokens in all 4,986 generated emails and none
of the 5,000 human emails. Such source-associated formatting
may affect parsed features; its contribution to performance
was not isolated.

Grouping detects similarity under a binary word-trigram
representation and the selected 0.88 cutoff. It cannot
guarantee that all paraphrases or related templates are
separated. Evaluation used one partition seed and 20 settings
per classifier. Shared test emails and overlapping fold
training sets create dependence; no inferential significance
tests or group-bootstrap confidence intervals were computed.

Recalibration requires labeled target data and optimizes F1
without an explicit false-positive cost constraint.
Pooling changes both generator diversity and training size,
and tests generators represented in training. Source holdout
also changes human training size and composition.

The SHAP investigation was exploratory and focused on CatBoost
with one paired LightGBM diagnostic. Across-study rankings
use different fitted models and explanation records, and
SHAP contributions do not demonstrate causal effects.
The experiments did not evaluate legitimate-email detection,
new generator families, future time periods, or operational
header, reputation, and attachment signals.

\section{Conclusion}
\label{sec:conclusion}

This study reused 9,986 phishing emails and reproduced their
17 released stylometric features to compare KNN, LightGBM,
CatBoost, XGBoost, and Random Forest. Python pipelines applied
shared partitions, group-aware validation, training-only KNN
scaling, and separate calibration data. Five evaluations
examined familiar generators, generator changes, threshold
adjustment, human-source holdout, and pooled training.

The strongest classifier varied by evaluation condition.
Recalibration improved mean transfer F1 for all five classifiers,
with a trade-off in human false positives, while pooling
achieved high performance on the represented generators.
CatBoost SHAP agreed with published XGBoost top-five membership
and added an exploratory diagnostic of missed Llama emails.
The accountability ledger preserves the evidence needed to
trace these findings and extend the experiments.

\subsection{Future Work}

Future studies should reserve fresh test data for independent
corpora, additional generators, alternative prompts, and
multilingual evaluations. Controlled formatting ablations
and equal-size training comparisons could separate preparation
effects from generator diversity. Repeated partitions and
group-aware uncertainty estimates would strengthen comparisons.
Further work should examine calibration sample size and
false-positive constraints, compare XGBoost and CatBoost SHAP
on identical explanation records, and evaluate legitimate emails
within a broader security pipeline. The ledger can retain
these future experiments without introducing automatic retraining.

% Finish pending result floats without forcing an extra bibliography page.
\FloatBarrier
% Let unused column space remain at the bottom, rather than stretching
% the gaps below the References heading or between bibliography entries.
\raggedbottom

\begingroup
\setlength{\emergencystretch}{2em}

\begin{thebibliography}{99}
\setlength{\itemsep}{0pt}
\setlength{\parsep}{0pt}

\bibitem{salloum2022}
S. A. Salloum, T. Gaber, S. Vadera, and K. Shaalan,
``A systematic literature review on phishing email detection using
natural language processing techniques,''
\emph{IEEE Access}, vol. 10, pp. 65703--65727, 2022.
doi: \url{https://doi.org/10.1109/ACCESS.2022.3183083}.

\bibitem{heiding2024}
F. Heiding, B. Schneier, A. Vishwanath, J. Bernstein,
and P. S. Park,
``Devising and detecting phishing emails using large language models,''
\emph{IEEE Access}, vol. 12, pp. 42131--42146, 2024.
doi: \url{https://doi.org/10.1109/ACCESS.2024.3375882}.

\bibitem{opara2025}
C. Opara, P. Modesti, and L. Golightly,
``Evaluating spam filters and stylometric detection of AI-generated
phishing emails,''
\emph{Expert Systems with Applications}, vol. 276,
Art. no. 127044, 2025.
doi: \url{https://doi.org/10.1016/j.eswa.2025.127044}.

\bibitem{koide2026}
T. Koide, N. Fukushi, H. Nakano, and D. Chiba,
``ChatSpamDetector: Leveraging large language models for effective
phishing email detection,''
in \emph{Security and Privacy in Communication Networks},
Lecture Notes of the Institute for Computer Sciences, Social Informatics
and Telecommunications Engineering, vol. 629.
Cham, Switzerland: Springer, 2026, pp. 297--319.
doi: \url{https://doi.org/10.1007/978-3-031-94455-0_14}.

\bibitem{li2024mage}
Y. Li, Q. Li, L. Cui, W. Bi, Z. Wang, L. Wang, L. Yang,
S. Shi, and Y. Zhang,
``MAGE: Machine-generated text detection in the wild,''
in \emph{Proc. 62nd Annual Meeting of the Association for
Computational Linguistics}, vol. 1, 2024, pp. 36--53.
doi: \url{https://doi.org/10.18653/v1/2024.acl-long.3}.

\bibitem{borile2025}
C. Borile and C. Abrate,
``How to generalize the detection of AI-generated text:
Confounding neurons,''
in \emph{Findings of the Association for Computational Linguistics:
EMNLP 2025}, 2025, pp. 25461--25476.
doi: \url{https://doi.org/10.18653/v1/2025.findings-emnlp.1388}.

\bibitem{gutierrez2026}
R. Gutierrez, W. Villegas-Ch, and J. Govea,
``Cross-model evaluation of phishing detectors against
LLM-generated emails,''
\emph{Frontiers in Big Data}, vol. 9,
Art. no. 1883452, 2026.
doi: \url{https://doi.org/10.3389/fdata.2026.1883452}.

\bibitem{chataut2024}
R. Chataut, P. K. Gyawali, and Y. Usman,
``Can AI keep you safe? A study of large language models for
phishing detection,''
in \emph{Proc. IEEE 14th Annual Computing and Communication
Workshop and Conference (CCWC)}, 2024, pp. 548--554.
doi: \url{https://doi.org/10.1109/CCWC60891.2024.10427626}.

\bibitem{opara2024styloai}
C. Opara,
``StyloAI: Distinguishing AI-generated content with stylometric analysis,''
in \emph{Artificial Intelligence in Education. Posters and Late Breaking
Results, Workshops and Tutorials, Industry and Innovation Tracks,
Practitioners, Doctoral Consortium and Blue Sky},
Communications in Computer and Information Science, vol. 2151.
Cham, Switzerland: Springer, 2024, pp. 105--114.
doi: \url{https://doi.org/10.1007/978-3-031-64312-5_13}.

\bibitem{gutierrezdataset2026}
R. Gutierrez, W. Villegas-Ch, and J. Govea,
``Cross-model evaluation of phishing detectors against
LLM-generated emails: dataset, code and results,''
Zenodo, ver. 1, May 2026.
doi: \url{https://doi.org/10.5281/zenodo.20250116}.

\bibitem{lundberg2017}
S. M. Lundberg and S.-I. Lee,
``A unified approach to interpreting model predictions,''
in \emph{Advances in Neural Information Processing Systems},
vol. 30, 2017.

% The following entries were retained from the supplied draft.
% They currently have no \cite in the manuscript body.
% Before submission, retain only sources actually supporting the manuscript.
\bibitem{apwg2025}
Anti-Phishing Working Group,
``Phishing activity trends report, 4th quarter 2025,''
APWG, Technical Report, 2025.

\bibitem{asiri2023}
S. Asiri, Y. Xiao, S. Alzahrani, S. Li, and T. Li,
``A survey of intelligent detection designs of HTML URL
phishing attacks,''
\emph{IEEE Access}, vol. 11, pp. 6421--6443, 2023.
doi: \url{https://doi.org/10.1109/ACCESS.2023.3237798}.

\bibitem{bountakas2021}
P. Bountakas, K. Koutroumpouchos, and C. Xenakis,
``A comparison of natural language processing and machine learning
methods for phishing email detection,''
in \emph{Proc. 16th International Conference on Availability,
Reliability and Security (ARES)}, 2021.
doi: \url{https://doi.org/10.1145/3465481.3469205}.

\bibitem{ceas2008}
CEAS,
``CEAS 2008 live spam challenge corpus,''
2008. [Online]. Available:
\url{https://www.kaggle.com/datasets/naserabdullahalam/phishing-email-dataset}

\bibitem{cormack2007}
G. V. Cormack,
``TREC 2007 spam track overview,''
in \emph{Proc. 16th Text Retrieval Conference (TREC)},
National Institute of Standards and Technology, 2007.
doi: \url{https://doi.org/10.6028/NIST.SP.500-274.spam-overview}.

\bibitem{gore2023}
S. Gore, S. Hamsa, S. Roychowdhury, G. Patil, S. Gore,
and S. Karmode,
``Augmented intelligence in machine learning for cybersecurity:
Enhancing threat detection and human-machine collaboration,''
in \emph{Proc. 2nd International Conference on Augmented
Intelligence and Sustainable Systems (ICAISS)}, 2023,
pp. 638--644.
doi: \url{https://doi.org/10.1109/ICAISS58487.2023.10250514}.

\bibitem{heiding2026}
F. Heiding, S. Lermen, A. Kao, C. Mayrink Verdun,
B. Schneier, and A. Vishwanath,
``Evaluating large language models' ability to automate
spear phishing,''
\emph{Expert Systems with Applications}, vol. 314,
Art. no. 131546, 2026.
doi: \url{https://doi.org/10.1016/j.eswa.2026.131546}.

\bibitem{klimt2004}
B. Klimt and Y. Yang,
``The Enron corpus: A new dataset for email classification research,''
in \emph{Machine Learning: ECML 2004}.
Berlin, Germany: Springer, 2004, pp. 217--226.
doi: \url{https://doi.org/10.1007/978-3-540-30115-8_22}.

\bibitem{nazario2007}
J. Nazario,
``The Nazario phishing corpus,''
2007. [Online]. Available:
\url{https://monkey.org/~jose/phishing/}

\bibitem{radev2008}
D. R. Radev,
``CLAIR collection of fraud email,''
2008. [Online]. Available:
\url{https://aclweb.org/aclwiki/CLAIR_collection_of_fraud_email_(Repository)}

\bibitem{roy2024}
S. S. Roy, P. Thota, K. V. Naragam, and S. Nilizadeh,
``From chatbots to phishbots? Phishing scam generation in
commercial large language models,''
in \emph{Proc. IEEE Symposium on Security and Privacy (SP)},
2024, pp. 36--54.
doi: \url{https://doi.org/10.1109/SP54263.2024.00182}.

\bibitem{sakkis2003}
G. Sakkis, I. Androutsopoulos, G. Paliouras, V. Karkaletsis,
C. D. Spyropoulos, and P. Stamatopoulos,
``A memory-based approach to anti-spam filtering for mailing lists,''
\emph{Information Retrieval}, vol. 6, pp. 49--73, 2003.
doi: \url{https://doi.org/10.1023/A:1022948414856}.

\bibitem{verizon2025}
Verizon,
``2025 data breach investigations report,''
Technical Report, 2025. [Online]. Available:
\url{https://www.verizon.com/business/resources/Tc25/reports/2025-dbir-data-breach-investigations-report.pdf}

\bibitem{xu2025}
H. Xu, S. Wang, N. Li, K. Wang, Y. Zhao, K. Chen,
\emph{et al.},
``Large language models for cyber security:
A systematic literature review,''
\emph{ACM Transactions on Software Engineering and Methodology},
2025.
doi: \url{https://doi.org/10.1145/3769676}.

\bibitem{zhang2025}
J. Zhang, H. Bu, H. Wen, Y. Liu, H. Fei, R. Xi,
\emph{et al.},
``When LLMs meet cybersecurity: A systematic literature review,''
\emph{Cybersecurity}, vol. 8, Art. no. 55, 2025.
doi: \url{https://doi.org/10.1186/s42400-025-00361-w}.

\end{thebibliography}
\endgroup

\end{document}
